% ch04_u07 —— 原书 p.182–p.191（4.5.11 的证明、4.5.12、4.6.1–4.6.20）
% 本单元自 p.182 起，开篇为命题 4.5.11 的证明（该命题的陈述在 p.181，属上一单元）。
\SourcePage{182}

\emph{证明.} 设 $x=gG(O)\in\mathcal{GR}$。利用 $\alpha^{\vee}$ 与伴随表示之间的关系（见 4.5.10），以及 4.3.4 的注 (ii) 与 (iii)，可见 $x$ 属于 $\mathcal{GR}$ 的偶分支，当且仅当
\begin{equation}\tag{231}
\dim\mathfrak{g}_{ss}\otimes O/((\mathfrak{g}_{ss}\otimes O)\cap\operatorname{Ad}_g(\mathfrak{g}_{ss}\otimes O))
\end{equation}
为偶数。而（231）正是 $x$ 的 $G(O)$-轨道的维数。\hfill$\Box$

另有如下证明。利用（229）与 4.5.9 的注，可见该命题等价于公式 $\chi_Z(\alpha(-1))=(-1)^{\langle\chi,2\rho\rangle}$。此式显然：按（56），$\alpha:\mu_2\to Z$ 是相应于 $2\rho$ 的态射 $\lambda^{\#}:\mathbb{G}_m\to H\subset G$ 的限制。

\bdp{4.5.12} 下文关于 $\mathcal{GR}$ 中 $G(O)$-轨道的性质，本文不会用到，但我们认为仍值得一提。

$\overline{\operatorname{Orb}}_\chi$ 是诸 $\operatorname{Orb}_{\chi'}$（$\chi'\le\chi$）的并。事实上，设 $\rho:G\to GL(V)$ 是以 $\lambda$ 为最低权的表示，则对 $g\in\operatorname{Orb}_\chi$ 有 $\rho(g)\in t^{(\chi,\lambda)}\operatorname{End}(V\otimes O)$，而 $\rho(g)\notin t^{(\chi,\lambda)+1}\operatorname{End}(V\otimes O)$。于是若 $\operatorname{Orb}_{\chi'}\subset\overline{\operatorname{Orb}}_\chi$，则对 $G$ 的每个反支配权 $\lambda$ 都有 $(\chi-\chi',\lambda)\le0$。因此 $\chi-\chi'$ 是 $G$ 的单余根的线性组合，系数非负；由 4.5.4(i)，这些系数是整数，故 $\chi'\le\chi$。另一方面，一次 $GL(2)$ 计算表明：${}^{L}G$ 中满足 $\operatorname{Orb}_{\chi'}\subset\overline{\operatorname{Orb}}_\chi$ 的权 $\chi'$ 构成的集合，在 [Bour75] 第 VIII 章 §7 第 2 号的意义下是饱和的。于是 loc.cit 的命题 5 表明：对每个满足 $\chi'\le\chi$ 的支配权 $\chi'$，都有 $\operatorname{Orb}_{\chi'}\subset\overline{\operatorname{Orb}}_\chi$。

上述对 $\overline{\operatorname{Orb}}_\chi$ 的描述表明：$\operatorname{Orb}_\chi$ 闭，当且仅当 $\chi$ 极小。若 $G$ 单，则 $\chi$ 极小当且仅当 $\chi=0$，或 $\chi$ 是 ${}^{L}G$ 的微权（见 [Bour68] 第 VI 章 §2 习题 5）。于是 $\mathcal{GR}$ 的每个连通分支上恰有一个闭的 $G(O)$-轨道（用 4.5.4 与 loc.cit 习题的第一部分）。若 $\operatorname{Orb}_\chi$ 闭，则它是射影的；此时 $G(O)$ 经 $G=G(O/tO)$ 作用在 $\operatorname{Orb}_\chi$ 上，而 $\operatorname{Orb}_\chi$ 是 $G$ 对某个抛物子群的商。用 9.1.3 的话说，$\operatorname{Orb}_\chi=\operatorname{orb}_\chi=G/P_\chi^-$。

\SourcePage{183}

设 $G$ 单。则恰有一个 $\chi$，使得 $\overline{\operatorname{Orb}}_\chi\setminus\operatorname{Orb}_\chi$ 由单个点构成\footnote{当然，该点是 $e\in G(K)$ 的像。}；这个 $\chi$ 是 $\mathfrak{g}:=\operatorname{Lie}G$ 中相应于 $\mathfrak{g}$ 的最大根 $\alpha_{max}$ 的余根（见 [Bour75] 第 VIII 章 §7 习题 22）。此时 $\overline{\operatorname{Orb}}_\chi$ 可描述如下。置 $V:=\mathfrak{g}\otimes(\mathfrak{m}^{-1}/O)$，其中 $\mathfrak{m}$ 是 $O$ 的极大理想。记 $\bar V$ 为以 $V$ 为仿射子空间的射影空间。于是 $\bar V$ 是 $V\oplus\mathbb{C}$ 中直线构成的空间；特别地，$V^*=\mathfrak{g}^*\otimes(\mathfrak{m}/\mathfrak{m}^2)$ 作用在 $\bar V$ 上，并保持 $0\in V$。记 $C$ 为 $V$ 中那些 $G$-共轭于 $\mathfrak{g}_{\alpha_{max}}\otimes(\mathfrak{m}^{-1}/O)$ 的元素构成的集合。这是 $V$ 的闭子簇。其射影闭包 $\bar C\subset\bar V$ 是 $V^*$-不变的，因为 $C$ 是锥。容易证明，态射 $\exp:C\to G(K)/G(O)$ 可延拓为同构 $f:\bar C\xrightarrow{\ \sim\ }\overline{\operatorname{Orb}}_\chi$。显然 $f$ 既是 $\operatorname{Aut}^0 O$-等变的，也是 $G$-等变的。$\operatorname{Ker}(G(O)\to G(O/\mathfrak{m}))$ 在 $\bar C$ 上的作用由它在 $\overline{\operatorname{Orb}}_\chi$ 上的作用诱导，该作用来自 $V^*$ 在 $\bar C$ 上的作用，以及同构
\[
\operatorname{Ker}(G(O/\mathfrak{m}^2)\to G(O/\mathfrak{m}))\xrightarrow{\ \sim\ }\mathfrak{g}\otimes\mathfrak{m}/\mathfrak{m}^2\xrightarrow{\ \sim\ }V^*,
\]
其中最后一个箭头由 $\mathfrak{g}$ 上使 $(\alpha_{max},\alpha_{max})=2$ 的不变内积诱导。

\bdsec{4.6. 局部 Pfaffian 丛}

考虑仿射 Grassmann 流形 $\mathcal{GR}:=G(K)/G(O)$，其中 $O=\mathbb{C}[[t]]$，$K=\mathbb{C}((t))$。置 $Z:=\operatorname{Hom}(\pi_1(G),\mathbb{G}_m)$（由 4.1.1 的注，$Z$ 是 ${}^{L}G$ 的中心）。本小节将构造并研究一个函子 $\mathcal{L}\mapsto\lambda_{\mathcal{L}}=\lambda_{\mathcal{L}}^{\mathrm{loc}}$，它从群胚 $Z\operatorname{tors}_\theta(O)$（见 3.4.3）指向 $\mathcal{GR}$ 上线丛的范畴。称 $\lambda_{\mathcal{L}}$ 为相应于 $\mathcal{L}$ 的\emph{局部 Pfaffian 丛}。

我们建议读者眼下跳过本小节。

\bdp{4.6.1} 在 4.4.9 中我们定义了函子 $\mathcal{L}\mapsto\widetilde{G(K)}_{\mathcal{L}}$，它从 $Z\operatorname{tors}_\theta(O)$ 指向 $G(K)$ 被 $\mathbb{G}_m$ 的中心扩张的范畴。对 $\mathcal{L}\in Z\operatorname{tors}_\theta(O)$ 有分裂 $G(O)\to\widetilde{G(K)}_{\mathcal{L}}$，因而有主 $\mathbb{G}_m$-丛
\begin{equation}\tag{232}
\widetilde{G(K)}_{\mathcal{L}}/G(O)\to G(K)/G(O)=\mathcal{GR}.
\end{equation}

\SourcePage{184}

\bdp{4.6.2} \emph{定义.} $\lambda_{\mathcal{L}}$ 是 $\mathcal{GR}$ 上相应于 $\mathbb{G}_m$-丛（232）的线丛之逆。

显然 $\lambda_{\mathcal{L}}$ 函子性地依赖于 $\mathcal{L}\in Z\operatorname{tors}_\theta(O)$。

\bdp{4.6.3} \emph{注.} $\widetilde{G(K)}_{\mathcal{L}}$ 依赖于 $\mathfrak{g}$ 上非退化不变双线性型的选取（见 4.4.7）。故 $\lambda_{\mathcal{L}}$ 亦然。

\bdp{4.6.4} 记 $\bar e\in\mathcal{GR}$ 为单位元 $e\in G$ 的像。我们的 $\lambda_{\mathcal{L}}$ 是 $\mathcal{GR}$ 上唯一的 $\widetilde{G(K)}_{\mathcal{L}}$-等变线丛：它在 $\bar e$ 上平凡化，且任意 $c\in\mathbb{G}_m\subset\widetilde{G(K)}_{\mathcal{L}}$ 通过乘以 $c^{-1}$ 作用在 $\lambda_{\mathcal{L}}$ 上。唯一性由等式 $\operatorname{Hom}(G(O),\mathbb{G}_m)=0$ 推出。

\bdp{4.6.5} 由 4.4.11，$\widetilde{G(K)}_{\mathcal{L}}$ 在 $\lambda_{\mathcal{L}}$ 上的作用诱导出 $\widetilde{\mathfrak{g}}\otimes K$ 在 $\lambda_{\mathcal{L}}$ 上的作用，使得 $1\in\mathbb{C}\subset\widetilde{\mathfrak{g}}\otimes K$ 通过乘以 $-1$ 作用。它与 $\mathfrak{g}\otimes K$ 经左无穷小平移在 $\mathcal{GR}$ 上的作用相容。

\bdp{4.6.6} （63）经态射（56）的前推是正合列
\begin{equation}\tag{233}
0\to Z\to\operatorname{Aut}_Z O\to\operatorname{Aut}O\to0.
\end{equation}
对任意 $\mathcal{L}\in Z\operatorname{tors}(O)$，正合列
\begin{equation}\tag{234}
0\to Z\to\operatorname{Aut}(O,\mathcal{L})\to\operatorname{Aut}O\to0
\end{equation}
可与（233）典范地等同。这里 $\operatorname{Aut}(O,\mathcal{L})$ 是由对 $(\sigma,\varphi)$ 构成的群归纳概形，其中 $\sigma\in\operatorname{Aut}O$，$\varphi:\mathcal{L}\xrightarrow{\ \sim\ }\sigma_*\mathcal{L}$（读者或许更愿意把 $\mathcal{L}$ 视为 3.4.5 中范畴 $\widetilde{Z}\operatorname{tors}_\omega(O)$ 的对象）。（233）与（234）之间的同构由显然的态射 $\operatorname{Aut}_2 O:=\operatorname{Aut}(O,\omega_O^{1/2})\to\operatorname{Aut}(O,\mathcal{L})$ 诱导。

$\operatorname{Aut}_Z O=\operatorname{Aut}(O,\mathcal{L})$ 通过结构的转移作用在正合列（217）上；$\operatorname{Aut}_Z O$ 在 $\mathbb{G}_m$ 上的作用平凡，其在 $G(K)$ 上的作用来自 $\operatorname{Aut}O$ 在 $G(K)$ 上的通常作用。子群 $G(O)\subset\widetilde{G(K)}_{\mathcal{L}}$ 是 $\operatorname{Aut}_Z O$-不变的。

\SourcePage{185}

\bdp{4.6.7} 由 4.6.6，$\operatorname{Aut}O$ 在 $\mathcal{GR}$ 上的作用典范地提升为 $\operatorname{Aut}_Z O$ 在主丛（232）与线丛 $\lambda_{\mathcal{L}}$ 上的作用。$\operatorname{Aut}_Z O$ 在 $\lambda_{\mathcal{L}}$ 上的作用诱导出 $\operatorname{Der}O=\operatorname{Lie}\operatorname{Aut}_Z O$ 在 $\lambda_{\mathcal{L}}$ 上的作用。

\bdp{4.6.8} $Z=\operatorname{Aut}\mathcal{L}$ 在扩张（217）上的作用来自（215）。故 $Z$ 经由态射
\begin{equation}\tag{235}
Z\to H^0(\mathcal{GR},\mathcal{O}^*_{\mathcal{GR}})
\end{equation}
作用在 $\lambda_{\mathcal{L}}$ 上，该态射是（215）与自然嵌入 $\operatorname{Hom}(G(K),\mathbb{G}_m)\hookrightarrow H^0(\mathcal{GR},\mathcal{O}^*_{\mathcal{GR}})$ 之复合的逆。回顾 $\pi_0(\mathcal{GR})=Z^{\vee}$（见 4.5.9），故 $z\in Z$ 定义了 $f_z:\pi_0(\mathcal{GR})\to\mathbb{C}^*$，而（235）是映射 $z\mapsto f_z^{-1}$。

\bdp{4.6.9} \emph{注.}（这里需要它吗 ???）把 $\mathcal{GR}$ 上线丛的范畴视为 3.4.4 意义下的 $Z$-范畴，$Z$-结构由（235）定义。由 3.4.7(i) 有典范的 Picard 函子
\begin{equation}\tag{236}
Z\operatorname{tors}(O)=Z\operatorname{tors}\to\{\text{$\mathcal{GR}$ 上的线丛}\}.
\end{equation}
确切地说，（236）把 $\mathcal{E}\in Z\operatorname{tors}$ 对应到 $\mathcal{O}_{\mathcal{GR}}$ 的 $\mathcal{E}$-扭，并配以 $Z$-作用（235）。由 3.4.7(iv)，函子 $\mathcal{L}\mapsto\lambda_{\mathcal{L}}$（$\mathcal{L}\in Z\operatorname{tors}_\theta(O)$）关于 Picard 函子（236）是仿射的。
% [Tutor Note: 原书 p.185 4.6.9 中的「???」是作者自留的行内疑问记号，不是排印故障；底本如此，保留原样。]

\bdp{4.6.10} 由（56）定义的态射 $\alpha:\mu_2\to Z$ 诱导出 $\mu_2$ 在 $\lambda_{\mathcal{L}}$ 上的作用（$\mathcal{L}\in Z\operatorname{tors}_\theta(O)$）。它给出 $\lambda_{\mathcal{L}}$ 上的 $(\mathbb{Z}/2\mathbb{Z})$-分次。在 4.5.10 中我们引入了 $\mathcal{GR}$ 的偶分支与奇分支的概念。由 4.6.8，$(\mathbb{Z}/2\mathbb{Z})$-分次丛 $\lambda_{\mathcal{L}}$ 到 $\mathcal{GR}$ 的偶（相应地，奇）分支上的限制是偶（相应地，奇）的。

\bdp{4.6.11} 函子
\begin{equation}\tag{237}
Z\operatorname{tors}_\theta(O)\to\{\text{$\mathcal{GR}$ 上的线丛}\},\qquad \mathcal{L}\mapsto\lambda_{\mathcal{L}}
\end{equation}
是 3.4.4 意义下的 $Z$-函子，只要（237）右端的 $Z$-结构由（235）定义。由于 $Z\operatorname{tors}_\theta(O)$ 等价于 $\omega^{1/2}(O)\otimes_{\mu_2}Z$\SourcePage{186}（见 3.4.4），函子（237）可由相应的函子
\begin{equation}\tag{238}
\omega^{1/2}(O)\to\{\text{$\mathcal{GR}$ 上的线丛}\}
\end{equation}
重构，其中 $\omega^{1/2}(O)$ 是 $\omega(O)$ 的平方根构成的群胚。由于扩张（212）本质上来自「Clifford 扩张」（193），给出（238）的 Clifford 式描述是容易的。答案如下。

设 $\mathcal{L}\in\omega^{1/2}(O)$。我们固定了 $\mathfrak{g}$ 上非退化不变对称双线性型，故 Tate 空间 $V=V_{\mathcal{L}}:=\mathcal{L}\otimes_O(\mathfrak{g}\otimes K)$ 带有非退化对称双线性型（见 4.3.3），而 $L:=\mathcal{L}\otimes\mathfrak{g}\subset V$ 是 Lagrange c-格。置 $M=M_{\mathcal{L}}:=\operatorname{Cl}(V)/\operatorname{Cl}(V)L$；这是 $\operatorname{Cl}(V)$ 上不可约的 $(\mathbb{Z}/2\mathbb{Z})$-分次离散模。在 $V$ 中 Lagrange c-格构成的归纳概形 $\operatorname{Lagr}(V)$ 上有线丛 $\mathcal{P}_M$（见 4.3.2）。我们断言
\begin{equation}\tag{239}
\lambda_{\mathcal{L}}=\varphi^*\mathcal{P}_{M_{\mathcal{L}}},
\end{equation}
其中态射\footnote{容易证明 $\varphi$ 是闭嵌入，其像是满足 $O\Lambda=\Lambda$ 且 $\mathcal{L}^{-1}\otimes_O\Lambda$ 为 $\mathfrak{g}\otimes K$ 的李子代数的那些 $\Lambda\in\operatorname{Lagr}(V)$ 构成的归纳概形。} $\varphi:G(K)/G(O)\to\operatorname{Lagr}(V)$ 由 $\varphi(g):=gLg^{-1}$ 定义；换言之，
\begin{equation}\tag{240}
\parbox{0.82\linewidth}{\raggedright $\lambda_{\mathcal{L}}$ 在 $g\in G(K)/G(O)$ 处的纤维为 $M^{gLg^{-1}}:=\{m\in M_{\mathcal{L}}\mid(gLg^{-1})\cdot m=0\}$。}
\end{equation}
事实上，中心扩张（212）与由（193）诱导的扩张相反，因此 $\widetilde{O}(V)$ 在 $\mathcal{P}_{M_{\mathcal{L}}}$ 上的作用（见 4.3.2）诱导出 $\widetilde{G(K)}_{\mathcal{L}}$ 在 $\varphi^*\mathcal{P}_{M_{\mathcal{L}}}$ 上的作用，使得 $c\in\mathbb{G}_m\subset\widetilde{G(K)}_{\mathcal{L}}$ 通过乘以 $c^{-1}$ 作用；此外，$\varphi^*\mathcal{P}_{M_{\mathcal{L}}}$ 在 $\bar e$ 上的纤维是 $\mathbb{C}$。

显然同构（239）关于 $\mathcal{L}\in\omega^{1/2}(O)$ 是函子性的。

\bdp{4.6.12} \emph{注.}
\begin{enumerate}[label=(\roman*),leftmargin=2.2em]
\item 4.3.2 中的线丛 $\mathcal{P}_M$ 是 $(\mathbb{Z}/2\mathbb{Z})$-分次的。故（239）两边都是 $(\mathbb{Z}/2\mathbb{Z})$-分次的。（239）两边的分次都由 $\mu_2=\operatorname{Aut}\mathcal{L}$ 的作用诱导（对右端论证这一点时注意：\SourcePage{187}$\operatorname{Cl}(V)$ 上的 $(\mathbb{Z}/2\mathbb{Z})$-分次由 $\mu_2$ 在 $V$ 上的自然作用诱导）。因此（239）是分次同构。
\item 由 4.6.10，$-1\in\mu_2=\operatorname{Aut}\mathcal{L}$ 通过乘以 $(-1)^p$ 作用在（239）的右端，其中 $p$ 是奇偶性函数（230）。这也可由等式 $\chi=\theta$（见引理 4.3.4 的证明）与 4.3.4 末的注 (ii) 推出。
\end{enumerate}

\bdp{4.6.13} 我们应当思考超的方面，尤其是：一维超空间的逆是什么？（或许应在任意 Picard 范畴中表述；若它不是\emph{严格}交换的，可能会出问题。）

考虑 $G(O)$-轨道 $\operatorname{Orb}_\chi\subset\mathcal{GR}$，其中 $\chi\in P_+({}^{L}G)$（见 4.5.8）。我们将计算 $\lambda_{\mathcal{L},\chi}:=\lambda_{\mathcal{L}}$ 到 $\operatorname{Orb}_\chi$ 的限制，$\mathcal{L}\in Z\operatorname{tors}_\theta(O)$。由 4.6.4，$\lambda_{\mathcal{L},\chi}$ 是 $G(O)$-等变的。轨道 $\operatorname{Orb}_\chi$ 是 $\operatorname{Aut}^0 O$-不变的；由 4.6.7，$\lambda_{\mathcal{L},\chi}$ 是 $\operatorname{Aut}^0_Z O$-等变的，其中 $\operatorname{Aut}^0_Z O$ 是 $\operatorname{Aut}^0 O$ 在 $\operatorname{Aut}_Z O$ 中的原像（见（233））。最后，$\lambda_{\mathcal{L},\chi}$ 是 $\mathbb{Z}/2\mathbb{Z}$-分次的（但实际上 $\lambda_{\mathcal{L},\chi}$ 依 $\chi$ 而为偶或奇；此外，$\mathbb{Z}/2\mathbb{Z}$-分次可由 $Z\subset\operatorname{Aut}^0_Z O$ 的作用重构……）。群 $G(O)$ 与 $\operatorname{Aut}^0_Z O$ 也作用在典范层 $\omega_{\operatorname{Orb}_\chi}$ 上（$\operatorname{Aut}^0_Z O$ 经 $\operatorname{Aut}^0 O$ 作用）。在 4.6.17--4.6.19（???）中，我们将对某个一维向量空间 $\mathfrak{d}_{\mathcal{L},\chi}$ 构造典范同构
\begin{equation}\tag{241}
\lambda_{\mathcal{L},\chi}\xrightarrow{\ \sim\ }\omega_{\operatorname{Orb}_\chi}\otimes(\mathfrak{d}_{\mathcal{L},\chi})^{-1}.
\end{equation}
该空间带有 $G(O)$ 与 $\operatorname{Aut}^0_Z O$ 的作用，且（241）关于这些群等变。
% [Tutor Note: 原书 p.187 4.6.13 中的「???」是作者自留的交叉引用疑问记号（针对 4.6.17--4.6.19 的段号）；底本如此，保留原样。]

\bdp{4.6.14} 我们来定义 $\mathfrak{d}_{\mathcal{L},\chi}$。当然，$G(O)$ 在 $\mathfrak{d}_{\mathcal{L},\chi}$ 上的作用规定为平凡（$G(O)$ 没有非平凡特征标）。所以须对每个 $\chi$ 构造函子
\begin{equation}\tag{242}
Z\operatorname{tors}_\theta(O)\to\{\operatorname{Aut}^0_Z O\text{-mod}\},\qquad \mathcal{L}\mapsto\mathfrak{d}_{\mathcal{L},\chi}
\end{equation}

\SourcePage{188}

其中 $\{\operatorname{Aut}^0_Z O\text{-mod}\}$ 表示 $\operatorname{Aut}^0_Z O$-模的范畴。先定义函子
\begin{equation}\tag{243}
\omega^{1/2}(O)\to\{\operatorname{Aut}^0_Z O\text{-mod}\},\qquad \mathcal{L}\mapsto\mathfrak{d}_{\mathcal{L},\chi}.
\end{equation}
对 $\mathcal{L}\in\omega^{1/2}(O)$ 置
\begin{equation}\tag{244}
\mathfrak{d}_{\mathcal{L},\chi}:=(\mathcal{L}_0)^{\otimes d(\chi)},
\end{equation}
其中 $\mathcal{L}_0$ 是 $\mathcal{L}$ 在闭点 $0\in\operatorname{Spec}O$ 上的纤维，而
\begin{equation}\tag{245}
d(\chi):=(\chi,2\rho)=\dim\operatorname{Orb}_\chi.
\end{equation}
如下定义 $\operatorname{Aut}^0_Z O$ 在 $\mathfrak{d}_{\mathcal{L},\chi}$ 中的表示：$\operatorname{Aut}^0_2 O=\operatorname{Aut}^0(O,\mathcal{L})$ 以显然的方式作用，而 $Z\subset\operatorname{Aut}^0_Z O$ 经由
\begin{equation}\tag{246}
\chi_Z:Z\to\mathbb{G}_m
\end{equation}
作用，其中 $\chi_Z$ 是 $\chi\in P^+({}^{L}G)$ 到 $Z\subset{}^{L}G$ 的限制（这两个作用相容，因为 $\chi_Z$ 与态射（56）的复合把 $-1\in\mu_2$ 映到 $(-1)^{(\chi,2\rho)}$）。

于是我们构造了（243）。$\omega^{1/2}(O)$ 是 3.4.4 意义下的 $\mu_2$-范畴，$\{\operatorname{Aut}^0_Z O\text{-mod}\}$ 是 $Z$-范畴，而（243）是 $\mu_2$-函子（$\{\operatorname{Aut}^0_Z O\}$ 上的 $\mu_2$-结构来自态射（56），等价地，来自典范嵌入 $\mu_2\to\operatorname{Aut}^0_2 O$）。故（243）诱导出 $Z$-函子
\[
Z\operatorname{tors}_\theta(O)=\omega^{1/2}(O)\otimes_{\mu_2}Z\to\{\operatorname{Aut}^0_Z O\text{-mod}\}.
\]
这就是（242）的定义。

\bdp{4.6.15} 显然 $\operatorname{Lie}\operatorname{Aut}^0_Z O=\operatorname{Der}^0 O$ 如下作用在一维空间 $\mathfrak{d}_{\mathcal{L},\chi}$ 上：
\begin{equation}\tag{247}
L_0\mapsto(\chi,\rho)=-\frac{1}{2}\dim\operatorname{Orb}_\chi,\qquad L_n\mapsto0\quad(n>0).
\end{equation}
按惯例，$L_n:=-t^{n+1}\frac{d}{dt}\in\operatorname{Der}^0 O$。
% [Tutor Note: 原书 p.188 把该特征标记作 $(\chi,\rho)$，而 (245) 已给出 $d(\chi)=(\chi,2\rho)=\dim\operatorname{Orb}_\chi$，两式相差一个符号。数值 $-\frac{1}{2}\dim\operatorname{Orb}_\chi$ 本身与约定 $L_n:=-t^{n+1}\frac{d}{dt}$ 相符（$L_0$ 在 $\omega^{1/2}$ 的纤维上乘 $-1/2$），故偏差只在记号 $(\chi,\rho)$。底本如此，保留原样。]

\SourcePage{189}

\bdp{4.6.16} \emph{注.} 4.6.14 中 $\mathfrak{d}_{\mathcal{L},\chi}$ 的定义可改述如下。利用 3.4.5 的等价 $Z\operatorname{tors}_\theta(O)\xrightarrow{\ \sim\ }\widetilde{Z}\operatorname{tors}_\omega(O)$，我们把 $\mathcal{L}\in Z\operatorname{tors}_\theta(O)$ 按（59）解释为 $\mathbb{G}_m$-torsor $\omega_O$ 到 $\widetilde{Z}$-torsor 的提升。由（62）有典范态射 $\widetilde{Z}\to{}^{L}H$，其中 ${}^{L}H$ 是 ${}^{L}G$ 的嘉当环面；或者同样地，${}^{L}H$ 是 ${}^{L}G$ 的一个嘉当子群，并配有一个包含它的固定 Borel 子群。记 $\chi_{\widetilde{Z}}$ 为 $\widetilde{Z}\to{}^{L}H$ 与 $\chi:{}^{L}H\to\mathbb{G}_m$ 的复合。$\operatorname{Spec}O$ 上的 $\widetilde{Z}$-torsor $\mathcal{L}$ 与一维表示 $\chi_{\widetilde{Z}}:\widetilde{Z}\to\mathbb{G}_m$ 定义了 $\operatorname{Spec}O$ 上的线丛 $\mathfrak{d}^{\mathcal{O}}_{\mathcal{L},\chi}$。按 4.6.6，$\operatorname{Aut}_Z O=\operatorname{Aut}(O,\mathcal{L})$；故 $\operatorname{Aut}O$ 在 $\operatorname{Spec}O$ 上的作用典范地提升为 $\operatorname{Aut}_Z O$ 在 $\mathfrak{d}^{\mathcal{O}}_{\mathcal{L},\chi}$ 上的作用。因此 $\operatorname{Aut}^0_Z O$ 作用在 $\mathfrak{d}^{\mathcal{O}}_{\mathcal{L},\chi}$ 在 $0\in\operatorname{Spec}O$ 处的纤维上。读者容易把该纤维与 4.6.14 中的 $\mathfrak{d}_{\mathcal{L},\chi}$ 等同。

\bdp{4.6.17} 我们来对 $\mathcal{L}\in\omega^{1/2}(O)$ 构造同构（241）。我们用 $\lambda_{\mathcal{L}}$ 的 Clifford 式描述。照 4.6.11 置 $V=V_{\mathcal{L}}:=\mathcal{L}\otimes_O(\mathfrak{g}\otimes K)$，$L:=\mathcal{L}\otimes\mathfrak{g}\subset V$，$M=M_{\mathcal{L}}:=\operatorname{Cl}(V)/\operatorname{Cl}(V)L$。对 $x\in\mathcal{GR}=G(K)/G(O)$ 置 $L_x:=gLg^{-1}$，其中 $g$ 是 $x$ 在 $G(K)$ 中的原像。由（240），$\lambda_{\mathcal{L}}$ 在 $x$ 处的纤维等于
\begin{equation}\tag{248}
M^{L_x}:=\{m\in M_{\mathcal{L}}\mid L_x\cdot m=0\}.
\end{equation}
设 $x\in\operatorname{Orb}_\chi$。因 $\operatorname{Orb}_\chi$ 是 $x$ 的 $G(O)$-轨道，$\operatorname{Orb}_\chi$ 在 $x$ 处的切空间为 $(\mathfrak{g}\otimes O)/((\mathfrak{g}\otimes O)\cap g(\mathfrak{g}\otimes O)g^{-1})=L^{-1}\otimes_O(L/(L\cap L_x))$，其中 $g\in G(K)$ 是 $x$ 的原像。于是 $\omega^{-1}_{\operatorname{Orb}_\chi}$ 在 $x$ 处的纤维等于 $(\mathcal{L}_0)^{\otimes-d(\chi)}\otimes\det({L/(L\cap L_x)})$，其中 $d(\chi)=\dim\operatorname{Orb}_\chi$。计及（244），可见（241）右端在 $x$ 处的纤维等于
\begin{equation}\tag{249}
(\det({L/(L\cap L_x)}))^{-1}.
\end{equation}
于是只需构造同构
\begin{equation}\tag{250}
\det({L/(L\cap L_x)})\otimes M^{L_x}\xrightarrow{\ \sim\ }\mathbb{C}.
\end{equation}
% [Tutor Note: 原书 p.189 用 $\omega^{-1}_{\operatorname{Orb}_\chi}$ 的纤维代入 (241)（(241) 本身写的是 $\omega_{\operatorname{Orb}_\chi}$）；再与 (244) 的 $\mathfrak{d}_{\mathcal{L},\chi}=(\mathcal{L}_0)^{\otimes d(\chi)}$ 联立，右端纤维为 $\det(L/(L\cap L_x))\otimes(\mathcal{L}_0)^{\otimes-2d(\chi)}$，与 (249) 所记的 $(\det(L/(L\cap L_x)))^{-1}$ 表面不符。疑原文 $\omega$ 与 $\omega^{-1}$ 的上标混用；底本如此，保留原样。]

\SourcePage{190}

\bdp{4.6.18} \emph{引理.} 考虑带非退化对称双线性型的 Tate 空间 $V$。设 $L,\Lambda\subset V$ 是 Lagrange c-格，$M$ 是 Clifford 代数 $\operatorname{Cl}(V)$ 上的不可约离散模。考虑由自然映射 $\bigwedge^dL\to\bigwedge^dV\to\operatorname{Cl}(V)$ 诱导的算子
\begin{equation}\tag{251}
\bigwedge\nolimits^{d}L\otimes M\to M.
\end{equation}
若 $d=\dim L/(L\cap\Lambda)$，则（251）诱导同构
\begin{equation}\tag{252}
\bigwedge\nolimits^{d}(L/(L\cap\Lambda))\otimes M^{\Lambda}\xrightarrow{\ \sim\ }M^{L}.
\end{equation}
证明可化归到 $\dim V<\infty$ 且 $V=L\oplus\Lambda$ 的情形。

\bdp{4.6.19} 我们定义（250）为 $\Lambda=L_x$ 时的同构（252）（在 4.6.17 的情形下 $M^L=\mathbb{C}$）。于是对 $\mathcal{L}\in\omega^{1/2}(O)$，我们构造了同构（241），它关于 $G(O)$ 与 $\operatorname{Aut}^0_2 O=\operatorname{Aut}^0(O,\mathcal{L})$ 等变。

记 $\mathcal{C}_\chi$ 为 $\operatorname{Orb}_\chi$ 上线丛的范畴。（241）两边都是 $\mu_2$-函子 $\omega^{1/2}(O)\to\mathcal{C}_\chi$，并延拓为 $Z$-函子
\[
Z\operatorname{tors}_\theta(O)=\omega^{1/2}(O)\otimes_{\mu_2}Z\to\mathcal{C}_\chi
\]
（$\mathcal{C}_\chi$ 上的 $Z$-结构由与（246）相逆的 $Z$ 的特征标定义）；对（241）左端，这由 4.6.8 推出。显然（241）是函子 $\omega^{1/2}(O)\to\mathcal{C}_\chi$ 的同构。因此（241）是函子 $Z\operatorname{tors}_\theta(O)\to\mathcal{C}_\chi$ 的同构。同构（241）是 $\operatorname{Aut}^0_Z O$-等变的，因为它既 $\operatorname{Aut}^0_2 O$-等变又 $Z$-等变。

\bdp{4.6.20} 回顾 $\lambda_{\mathcal{L}}$ 依赖于 $\mathfrak{g}$ 上非退化不变双线性型的选取（见 4.6.3 与 4.4.7）。如 4.4.7 的脚注所解释，$\lambda_{\mathcal{L}}$ 有更典范的版本。当 $G$ 单时，该版本 $\lambda_{\mathcal{L}}^{can}$ 依赖于 $\beta^{1/2}$ 的选取，其中 $\beta$ 是 $\mathfrak{g}$ 上不变双线性型构成的线（参看 4.4.5）；$\lambda_{\mathcal{L}}^{can}$ 来自下述版本的（212）：\SourcePage{191}%
以 $SO(\mathfrak{g}\otimes\beta^{1/2})$ 代替 $SO(\mathfrak{g})$。容易看出，$\lambda_{\mathcal{L}}^{can}$ 上相应于 $-1\in\operatorname{Aut}\beta^{1/2}$ 的作用的 $(\mathbb{Z}/2\mathbb{Z})$-分次，与 4.6.10 的分次一致。（241）的「典范」版本是同构
\begin{equation}\tag{253}
\lambda^{can}_{\mathcal{L},\chi}\xrightarrow{\ \sim\ }\omega_{\operatorname{Orb}_\chi}\otimes(\mathfrak{d}_{\mathcal{L},\chi})^{-1}\otimes(\beta^{1/2})^{\otimes-d(\chi)},
\end{equation}
其中 $d(\chi)$ 由（245）定义。细节留给读者。
